\chapter[\emj{⏱️}~Big-O Notation]{\emj{⏱️}~Big-O Notation}

\begin{dsaquees}

Imagínate que tienes una caja de chocolates 🍫 y quieres
encontrar uno de sabor fresa. ¿Cuánto tiempo tardarías?

\begin{itemize}
    \item \textbf{Caso 1:} Si sabes que el de fresa SIEMPRE está en la posición 1
          $\rightarrow$ Lo agarras directo. No importa si hay 10 o 1,000,000 chocolates.
          $\rightarrow$ Siempre es 1 paso. ¡Rápidísimo! ⚡

    \item \textbf{Caso 2:} Si los chocolates están desordenados y tienes que revisar
          uno por uno hasta encontrar el de fresa
          $\rightarrow$ Si hay 10 chocolates, revisas hasta 10.
          $\rightarrow$ Si hay 1,000,000 chocolates, revisas hasta 1,000,000.
          $\rightarrow$ El tiempo crece igual que la cantidad de chocolates. 📈

    \item \textbf{Caso 3:} Si por cada chocolate que revisas, tienes que revisar
          TODOS los demás para compararlos
          $\rightarrow$ Si hay 10 chocolates $\rightarrow$ 10 $\times$ 10 = 100 comparaciones
          $\rightarrow$ Si hay 100 chocolates $\rightarrow$ 100 $\times$ 100 = 10,000 comparaciones
          $\rightarrow$ ¡Crece muy rápido y se pone muy lento muy pronto! 🐌
\end{itemize}

Big-O Notation es la forma que tienen los programadores para
describir qué tan rápido o lento crece el tiempo que tarda
un algoritmo cuando le das más y más datos.

La "n" siempre representa la cantidad de datos de entrada.
Big-O siempre describe el PEOR caso posible (lo más lento que puede ser).

\end{dsaquees}

\begin{dsaotraforma}

Big-O es la métrica estándar para comparar la eficiencia de
algoritmos independientemente del hardware. Describe cómo
escala el tiempo de ejecución (o uso de memoria) en función
del tamaño del input "n", siempre en el worst case.

En interviews, SIEMPRE tienes que dar el Big-O de tu solución.
Es tan esperado como que el código funcione.

\end{dsaotraforma}

\begin{dsadiagrama}

\begin{verbatim}
Tiempo
  |                                              O(2^n) 💀
  |                                         ,,''
  |                                    ,,'''        O(n^2) 🐌
  |                               ,,'''        ,,,''
  |                          ,,'''       ,,,'''
  |                     ,,'''      ,,,'''          O(n log n) 🐢
  |                ,,'''      ,,'''
  |           ,,'''      ,,'''                    O(n) 🚶
  |      ,,'''       ,,''
  |  ,,''        ,,''                           O(log n) 🏃
  |''        ,,''
  |      ,,''                                  O(1) ⚡ (línea plana)
  +-------------------------------------------------> n (tamaño del input)
        pequeño      mediano        grande

  ⚡ O(1)       → Constante    → IDEAL
  🏃 O(log n)   → Logarítmico → MUY BUENO
  🚶 O(n)       → Lineal       → ACEPTABLE
  🐢 O(n log n) → Linealítmico → BUENO para sorting
  🐌 O(n^2)      → Cuadrático   → MALO (evitar si puedes)
  💀 O(2^n)     → Exponencial  → FATAL (solo para n muy pequeño)
\end{verbatim}

\end{dsadiagrama}

\begin{dsacomofunciona}

\begin{description}
    \item[O(1) --- Tiempo Constante] Sin importar cuántos datos tengas, siempre tarda lo mismo. Ejemplo: acceder a \texttt{arr[5]} directamente por índice.
    \item[O(log n) --- Tiempo Logarítmico] Cada paso DIVIDE el problema a la mitad. Si tienes 1,024 elementos $\rightarrow$ necesitas máximo 10 pasos ($2^{10} = 1024$). Ejemplo: Binary search.
    \item[O(n) --- Tiempo Lineal] Si duplicas el input, el tiempo se duplica. Ejemplo: recorrer todos los elementos de una lista una vez.
    \item[O(n log n) --- Tiempo Linealítmico] Un poco peor que lineal. La mayoría de buenos algoritmos de sorting. Ejemplo: Merge Sort.
    \item[$O(n^2)$ --- Tiempo Cuadrático] Dos loops anidados sobre el mismo input. Si duplicas el input, el tiempo se CUADRUPLICA. Ejemplo: Bubble Sort.
    \item[$O(2^n)$ --- Tiempo Exponencial] Cada elemento nuevo DUPLICA el trabajo. n=30 $\rightarrow$ más de 1,000,000,000 operaciones. Ejemplo: Fibonacci sin memoization.
\end{description}

\end{dsacomofunciona}

\begin{lstlisting}[language=C++]
#include <iostream>
#include <vector>
using namespace std;

// O(1) — Constante
int obtenerPrimero(vector<int>& arr) {
    return arr[0];  // acceso directo por índice → O(1)
}

// O(n) — Lineal
bool buscarLineal(vector<int>& arr, int objetivo) {
    for (int i = 0; i < arr.size(); i++) {
        if (arr[i] == objetivo) return true;
    }
    return false;
}

// O(n^2) — Cuadrático
bool tieneDuplicados(vector<int>& arr) {
    int n = arr.size();
    for (int i = 0; i < n; i++) {
        for (int j = i + 1; j < n; j++) {
            if (arr[i] == arr[j]) return true;
        }
    }
    return false;
}

// O(log n) — Logarítmico
int busquedaBinaria(vector<int>& arr, int objetivo) {
    int izq = 0;
    int der = arr.size() - 1;
    while (izq <= der) {
        int medio = izq + (der - izq) / 2;
        if (arr[medio] == objetivo) return medio;
        if (arr[medio] < objetivo)  izq = medio + 1;
        else                        der = medio - 1;
    }
    return -1;
}
\end{lstlisting}

\begin{dsacomplejidad}

\begin{tabular}{|l|l|l|l|}
\hline
Notación \& Nombre \& Ejemplo real \& ¿Usable? \\
\hline
O(1) \& Constante \& Array access, hash lookup \& ✅ Ideal \\
O(log n) \& Logarítmico \& Binary search, BST ops \& ✅ Muy bueno \\
O(n) \& Lineal \& Loop simple, linear search \& ✅ Aceptable \\
O(n log n) \& Linealítmico \& Merge sort, Heap sort \& ✅ Bueno \\
$O(n^2)$ \& Cuadrático \& Bubble sort, nested loops \& ⚠️ Malo \\
$O(2^n)$ \& Exponencial \& Subsets, recursion fatal \& ❌ Terrible \\
\hline
\end{tabular}

\end{dsacomplejidad}

\begin{dsaentrevistas}

\begin{itemize}
    \item Big-O ignora constantes y términos menores: $O(3n + 5) \rightarrow O(n)$.
    \item Siempre menciona tiempo Y espacio.
    \item \textbf{Regla de oro:} Si ves que tu solución tiene $O(n^2)$ y el problema involucra búsqueda, piensa si un hash map te da $O(n)$ en su lugar.
\end{itemize}

\end{dsaentrevistas}

\begin{dsaresumen}

\begin{itemize}
    \item Big-O = qué tan rápido crece el tiempo al crecer el input.
    \item Siempre describe el PEOR caso.
    \item $O(1) < O(\log n) < O(n) < O(n \log n) < O(n^2) < O(2^n)$.
    \item Hash maps convierten muchos $O(n^2)$ en $O(n)$.
\end{itemize}

\end{dsaresumen}


